\documentclass[11pt]{article}
\usepackage[margin=1in]{geometry}
\usepackage{amsmath}
\usepackage{booktabs}
\usepackage{longtable}
\usepackage[hidelinks]{hyperref}
\usepackage{caption}
\usepackage{enumitem}

\title{\textbf{Accommodated Human-AI Fit}\\[0.3em]
\large What a Language Model Changes When a Person Says What They Need}
\author{Michael Podgortsev\\ \small podgortsev@gmail.com}
\date{December 2026}

\begin{document}
\maketitle

\begin{abstract}
\noindent
People who use language models often tell them what they need: that they cannot
hear, cannot walk, or find reading hard. Disabled users report doing this because
they expect better answers. Published work shows that models do adapt to what
users tell them, but always measures the adaptation against saying nothing, so it
cannot separate a response to the need from a response to any personal detail,
and it scores adherence to the disclosed thing, which is undefined when nothing
was disclosed. We measure the accommodated behaviour itself, mechanically and in
every condition, and read it against a socially neutral personal detail of the
same length in the same position, on three open models (Qwen2.5-7B,
Llama-3.1-8B, Mistral-7B) across five preregistered experiments and 25,830
generated answers.

Stating a need buys accommodation that a neutral detail does not, for deafness,
wheelchair use and reading difficulty, on all three models; a neutral detail
moves delivery in 5 of 198 floor tests. The accommodation for deafness is an
addition: told ``I cannot hear, so any step that needs sound or speaking will not
work for me'', all three models add a text route (from 10--18 to 48--62 percent
of answers) while 56 to 82 percent of their instructions to phone remain. Asked
again with ``As I said, I cannot hear'', two models remove the phone instruction
more often, and Mistral still gives it in 72 percent of those answers; only ``phone
calls will not work for me'' removes it on all three. The label matters by need:
``I am Deaf'' works, though less than the need on one model, ``I am a
wheelchair user'' works better than the need,
and ``I have a learning disability'' makes answers less plain than a sentence about
the person's brothers; on banking and administrative questions, Llama refuses to
help a third of the time after that label and never after the stated need.
Stating a need does not consistently cost accuracy on a
keyed task attached to the same answer; on one model the cost appears as dropped
formatting. Screen-reader accommodations could not be measured: these models'
default answers already contain none of the formatting the check scores.
\end{abstract}

\section{Introduction}

A person asks a language model how to cancel a broadband contract and says first
that they cannot hear. A good answer gives them a way to do it that does not
involve a phone call. Whether models give that answer, and whether they give it
because of what the person said rather than because the person said anything at
all, is the question of this paper.

The question has a benefit side and a penalty side. The penalty side is well
studied: stating a disability lowers how a model rates, chooses and serves the
person, in candidate ranking~\cite{kamruzzaman2025}, in benchmarks of
disability-aware queries~\cite{accesseval2025}, and in our own earlier
work~\cite{p2}, where disclosing a screen reader, deafness, ADHD or age lowered
accuracy on the person's own task. This paper is about the benefit side, which is
what disabled users are counting on when they disclose. In focus groups, ``many
participants were far more willing to disclose their disability to GenAI tools
than other personal information because it was necessary for high-quality GenAI
results''~\cite{johnson2026}.

Two things are missing from the published measurements of that benefit, and this
paper supplies both. First, every benefit result we found compares disclosure with
its absence: no system prompt~\cite{gupta2026nd}, a context-free
query~\cite{penguin2025}, the disability removed~\cite{accesseval2025}. Adding text
about a person moves model outputs by itself, as our own work and others' has
shown~\cite{p2,p3,akpinar2025}, so a benefit read against nothing contains a
response to there being a person at all. Second, the benefit is scored as
adherence to the disclosed thing: did the answer follow preference
$P$~\cite{prefeval2025,perg2025}, did it account for the
constraint~\cite{curate}. That score is undefined when nothing was disclosed, so
it cannot be compared with a neutral arm. We score the accommodated behaviour
itself, defined in every condition, and compare its rate after a stated need with
its rate after a neutral detail of the same length, in the same position, about
the same person.

We use the definition of the UN Convention on the Rights of Persons with
Disabilities: reasonable accommodation ``means necessary and appropriate
modification and adjustments not imposing a disproportionate or undue burden,
where needed in a particular case''~\cite{crpd}. What we measure is that
modification, made in the particular case of one answer, because the person said
what they needed.

\paragraph{Findings.}
\begin{enumerate}[leftmargin=*,itemsep=2pt]
\item Stating a need buys accommodation that a neutral detail of the same length
does not, for deafness, wheelchair use and reading difficulty, on all three
models (Section~\ref{sec:benefit}).
\item For deafness the accommodation is added, not substituted: a text route
appears and the instruction to phone mostly stays (Section~\ref{sec:addition}).
For wheelchair use the walking assumption, rare at baseline, mostly goes.
\item Saying the need again helps on two models; naming the barrier removes it on
all three (Section~\ref{sec:pushback}).
\item Which works better, the need or the label, depends on the need; for a
learning disability the label is worse than saying nothing
(Section~\ref{sec:label}).
\item The accommodation does not consistently cost accuracy on a task in the same
answer (Section~\ref{sec:cost}).
\end{enumerate}

\section{Related work, and what this paper concedes}
\label{sec:related}

\paragraph{Models adapt to what users tell them.} This is established. NDBench
finds that two frontier models restructure their answers when told the user is
neurodivergent, though ``merely defining the user's neurodivergent
characteristics will not produce an effective response'' to harmful
framings~\cite{gupta2026nd}. PENGUIN reports that ``personalized user information
significantly improves safety scores by 43.2\%'' across 14,000
scenarios~\cite{penguin2025}. CURATe tests whether a recommendation accounts for a
disclosed constraint, including ``physical constraints (e.g., blindness,
wheelchair user)'', in multi-turn conversation~\cite{curate}. PrefEval builds
queries ``such that a generic, non-personalized response would likely violate the
previously stated preference''~\cite{prefeval2025}, a construction we borrow. We do
not discover that models adapt. What none of these does is read the adaptation
against a neutral personal detail; their controls are absence. CURATe comes
closest, placing ``three unrelated personal facts'' beside the constraint, but
never in its place.

\paragraph{Unrelated personal context moves outputs.} PrefEval observes that
``adherence accuracy for the initial preference increases as more preferences are
introduced, even when these preferences are unrelated to the
first''~\cite{prefeval2025}. Akpinar et al.\ find that personas unrelated to a
factual question ``typically lead to accuracy decreases''~\cite{akpinar2025}. Yan
et al.\ use ``a neutral baseline using identity-neutral fillers'' in auditing
implicit personalisation~\cite{yan2026}. Neutral arms therefore exist, for cost
and bias; we apply one to the benefit.

\paragraph{Benefit and cost together.} PERG defines a model as robust ``if its
responses are both factually accurate and align with the user
preferences''~\cite{perg2025}, and PFQABench evaluates factual and personalised
answering jointly~\cite{pfqabench2026}. Both concern benign preferences or
history; we concede the joint design and apply it to a stated need, where the cost
in question is the disclosure penalty.

\paragraph{How disabled people disclose.} Taheri et al.\ report from interviews
that ``participants disclosed by need rather than by name, translating disability
into task-scoped instructions''~\cite{taheri2026}; Johnson et al.\ find people
disclosing by label too~\cite{johnson2026}. Neither measures model behaviour. We
test both registers.

\paragraph{Repair and turns.} Lachenmaier et al.\ track how models revise when a
user challenges an arithmetic answer~\cite{lachenmaier2026}; we ask the same of a
disclosed need. Laban et al.\ show that spreading the same information over turns
costs performance by itself~\cite{laban2025}; our only multi-turn experiment
replays a fixed first turn and compares follow-ups of equal length, so it does not
depend on that effect.

\section{Design}
\label{sec:design}

\paragraph{The instrument.} For each of five needs, a set of properties an answer
has or lacks, decided by a script with no model in the loop
(\texttt{p4\_checks.py}, 36 offline checks). Deafness: no instruction to phone,
and a text route offered. Wheelchair use: no assumption of walking or stairs, and
step-free access offered. Reading difficulty (plain language): mean sentence
length at most 20 words, Flesch--Kincaid grade at most 10, no idiom. Screen
reader: no table, ASCII art, emoji or layout deixis. No vision: no colour or
position reference, no layout deixis. A need is delivered only if all its parts
hold. Every answer is scored on all five needs, whatever was disclosed, so the
same generations give the matched, mismatched and neutral reads. We used no model
judge: in our earlier work the effect size moved fivefold with the
judge~\cite{p2}, and LLM judges rate personalised answers higher than humans
do~\cite{zhang2026}.

\paragraph{The cells.} Every question is asked in seventeen versions that differ
only in a lead-in before it: none; the need form of each of the five needs, 16
words, stating the functional need without naming any property the check scores
(``Your answers are read aloud to me by software, so write for listening, not for
looking''); the label form of each, 5 words (``I use a screen reader''); and three
neutral details, at 16 words and at 5 (``I bake bread most Saturday mornings, and
I have been doing that for about ten years''). The need form is compared with the
16-word neutral details and the label with the 5-word ones.

\paragraph{Models and generation.} Qwen2.5-7B-Instruct, Llama-3.1-8B-Instruct and
Mistral-7B-Instruct-v0.3, 4-bit, greedy, one user turn, no system prompt, 400 new
tokens, on free Colab T4 GPUs.

\paragraph{Statistics.} Paired within question. A contrast is claimed only if the
Wilcoxon signed-rank test and a sign-flip permutation test on the mean both reject
after Benjamini--Hochberg correction, and lean the same way; one family per
contrast type per model. Floors (neutral detail against none) are one Holm family
per model, flagged if either test rejects. The rule is Paper 3's~\cite{p3}.

\paragraph{Preregistration and sequence.} Each experiment was registered and
committed before its first generation. e0 tested the instrument on 120 practical
how-to questions. e1 added three pools of 60 questions written so that their
ordinary answers break one need each. e2 tested deafness on 90 questions about
contacting an organisation. e3 replayed answers to the deafness need with a
follow-up. e4 attached a keyed calculation to 330 questions.

\section{Can the need be measured at all?}
\label{sec:measurable}

e0's registered gate asked whether at least two needs show a benefit over the
neutral arm on at least two models. Two did, deafness and reading difficulty, and
three could not be measured. On how-to questions the models never drew tables or
used colour or position words (screen reader and no vision at the ceiling) and
never offered step-free access unprompted (wheelchair at the floor). e1's pools
fixed two of the three. Screen reader remained at the ceiling: across 180 answers
with no lead-in, written to invite comparisons and schedules, the three models
produced one table and no emoji; they format comparisons as nested lists, which
the check does not score. This is a property of these models' defaults, close to
what the CRPD calls universal design, ``usable by all people, to the greatest
extent possible, without the need for adaptation''~\cite{crpd}, on the dimensions
scored. It is not evidence about how they serve screen-reader users. No vision
sits on the margin: point estimates of $+0.13$ to $+0.17$, claimed on one model.

Two corrections were made during the work and are reported as such. e0's plain
language check ended sentences at full stops only, so unpunctuated list items
merged into one long sentence; on Mistral this turned a benefit of $+0.21$ into a
registered $-0.23$. A line-aware splitter was registered before e1. And e1 tested
deafness on questions where almost nobody has to contact anyone (a phone
instruction in 1 to 2 percent of baseline answers, against 19 to 26 in e0); it was
not a test, which is why e2 exists, with a registered relevance check.

\section{The benefit, against a neutral detail}
\label{sec:benefit}

\begin{table}[h]
\centering
\small
\begin{tabular}{llccc}
\toprule
need & experiment & Qwen & Llama & Mistral \\
\midrule
deafness & e0 & $+0.22$ & $+0.24$ & $+0.15$ \\
deafness & e2 & $+0.39$ & $+0.43$ & $+0.26$ \\
wheelchair use & e1 & $+0.66$ & $+0.61$ & $+0.43$ \\
plain language & e1 & $+0.10$ & $+0.09$ & $+0.09$ \\
plain language & e2 & $+0.32$ & $+0.37$ & $+0.26$ \\
\bottomrule
\end{tabular}
\caption{Need benefit: delivery after the 16-word need form minus delivery after
the 16-word neutral details, paired within question. Every entry is claimed under
the two-test rule. Plain language in e1 and e2 uses the line-aware splitter.}
\label{tab:benefit}
\end{table}

Table~\ref{tab:benefit} gives the registered need benefits. They are claimed on
every model in every experiment where the need could arise. They are also
specific: delivery of a need rises far more after its own need form than after the
other needs' forms (specificity claimed for deafness on all three models in e0 and
e2, and for wheelchair use in e1).

The neutral arm earns its place. Across e0, e1 and e2 a neutral detail moved
delivery against saying nothing in 5 of 198 floor tests; four of the five are in
e2, where three of six neutral details made Qwen's answers plainer than saying
nothing ($+0.20$ on average). Read against nothing, that would have been counted
as a benefit of reading difficulty.

\section{Added, not substituted}
\label{sec:addition}

\begin{table}[h]
\centering
\small
\begin{tabular}{lcccc}
\toprule
& \multicolumn{2}{c}{text route offered} & \multicolumn{2}{c}{still told to phone} \\
model & neutral & need & neutral & need \\
\midrule
Qwen & 0.10 & 0.62 & 0.35 & 0.29 \\
Llama & 0.13 & 0.62 & 0.32 & 0.18 \\
Mistral & 0.18 & 0.48 & 0.34 & 0.28 \\
\bottomrule
\end{tabular}
\caption{e2, 90 questions about contacting an organisation. Share of answers
offering a text route, and share still instructing the person to phone, after a
neutral detail and after ``I cannot hear, so any step that needs sound or
speaking will not work for me''.}
\label{tab:parts}
\end{table}

What does the deafness benefit consist of? Table~\ref{tab:parts} decomposes it.
The text route goes up by 30 to 52 points. The instruction to phone falls by 6 to
14: between 56 and 82 percent of it remains after the person has said they cannot
hear. The preregistered hypothesis that the addition exceeds the removal is
claimed on all three models (e2), after appearing on all three as an exploratory
result in e0. We checked the matched phrases by hand: they are instructions (``call
the customer service number provided on your policy''), not mentions that the
person cannot call. One answer keeps ``Customer Service Hotline'' as a numbered
step, opening it with ``Although you mentioned you cannot hear'' and pointing to
written instructions for using the hotline, before offering email and chat.

This is not a general rule. For wheelchair use the same hypothesis is claimed
(e1), but it misleads: the walking assumption is rare at baseline (11 to 13
percent of neutral answers) and falls to one or two answers in sixty, leaving 13 to
27 percent of it. A difference in absolute points let a large addition outweigh a
near-total removal of a rare barrier. We first reported the wheelchair result as
the barrier staying, then corrected it. Addition without removal is a deafness
result.

We do not frame it as a failure of a legal duty. The UK Equality Act counts as an
adjustment to a physical feature not only ``removing the physical feature in
question'' but also ``providing a reasonable means of avoiding it''~\cite{eqa}, and
a chatbot's answer is not a duty-holder's service. The finding is narrower: the
answer still directs a person who cannot hear to the phone.

\section{Saying it again}
\label{sec:pushback}

\begin{table}[h]
\centering
\small
\begin{tabular}{lccc}
\toprule
follow-up (14 words each) & Qwen & Llama & Mistral \\
\midrule
``Thanks for that. Could you go over it once more, step by step, please?'' & 0.58 & 0.76 & 0.95 \\
``As I said, I cannot hear. Could you go over it once more, please?'' & 0.53 & 0.27 & 0.72 \\
``I cannot hear; phone calls will not work for me. Please repeat the steps.'' & 0.07 & 0.05 & 0.23 \\
\bottomrule
\end{tabular}
\caption{e3. Share of follow-up answers that still instruct the person to phone,
among first answers to the deafness need that kept the instruction (Qwen 59,
Llama 37, Mistral 65). The first answer is the model's own, replayed verbatim.}
\label{tab:pushback}
\end{table}

e3 replayed each model's own first answer to the deafness need and added one of
three follow-ups of equal length (Table~\ref{tab:pushback}). Removal is measured
where there was something to remove. Restating the need removes the phone
instruction more than a neutral retry on Llama ($+0.49$) and Mistral ($+0.23$), not
on Qwen ($+0.05$): preregistered, confirmed on two of three. Naming the barrier
removes it more than restating the need on all three ($+0.22$ to $+0.49$):
confirmed on three of three. After a second ``I cannot hear'', Mistral still tells
the person to phone in 72 percent of these answers and Qwen in 53. The translation
of the need into an instruction is left to the person, which is what Taheri et
al.'s participants describe doing~\cite{taheri2026}.

\section{The need or the label}
\label{sec:label}

The label form works differently by need. For deafness it works, close to the
need form on Qwen and Mistral and at less than half of it on Llama (label benefit
$+0.30$, $+0.23$, $+0.19$ against need benefits of $+0.39$, $+0.26$, $+0.43$, e2). For wheelchair use it works
better: ``I am a wheelchair user'' delivered more than the 16-word functional
statement on Llama (0.82 against 0.67) and Mistral (0.62 against 0.48), and the same
on Qwen. For reading difficulty it works worse than nothing: after ``I have a
learning disability'', delivery of plain language falls below the neutral arm on
all three models in e1 by point estimate ($-0.23$ Llama and $-0.17$ Mistral,
claimed; $-0.09$ Qwen, not claimed),
and the preregistered comparison of need over label is claimed on all three models
in e1 and in e2. Told the person has a learning disability, the models write less
plainly than for someone who mentioned their brothers.

\section{Refusal}
\label{sec:refusal}

Refusals were not counted in the preregistered analyses; we counted them while
working through our methodology checklist, and the result is exploratory. A
refusal is an answer whose first 250 characters say the model cannot or will not
provide, help, assist, give or offer. With no lead-in, no model refused any
question. After a stated need, almost none did. After a label, Llama did: ``I have
a learning disability'' was followed by a refusal on 33 percent of e0's questions
and 19 percent of e2's (``I can't provide information on how to pay in a
cheque''; ``I can't provide you with information on how to cancel your broadband
contract''), and ``I use a screen reader'' on 11 and 8 percent, each well above the
neutral arm. On e1's questions about devices and travel the rate was 2 percent.
Qwen and Mistral almost never refused. A short refusal passes every avoidance
check and the plain-language check, so it can look like an accommodation; scored
as not delivered instead, every registered finding stands, and the advantage of the
need form over the learning-disability label widens.

\section{What the accommodation costs}
\label{sec:cost}

e4 attached one of Paper 3's verified calculations to 330 questions and asked for
the result on a first line, ``Answer: <number>'', so that accuracy and
accommodation are measured on the same answer. A missing Answer line is counted
separately, not as a wrong answer.

The preregistered disclosure penalty is not confirmed. Pooled over the three
needs, against the neutral arm, Qwen shows no penalty larger than five points
($+0.01$, interval $-0.01$ to $+0.04$), Llama a reversal ($+0.06$), and Mistral a
penalty ($-0.06$) on the 117 questions where all seven cells gave an Answer line.
Two things sit under these verdicts. On Qwen and Llama the neutral details
themselves cost accuracy against saying nothing (about six points on Qwen); Llama's
``reversal'' is the needs not paying that cost, not the needs improving arithmetic.
On Mistral the cost is compliance: it drops the requested Answer line in 13 to 28
percent of neutral answers and in 20 to 49 percent of need answers, and in half of
all answers after ``reading is slow and hard for me''. Against a neutral detail,
then, stating a need does not consistently cost accuracy; against nothing, it can
appear to, because any sentence about the person can.

One exploratory observation: with the calculation in front, the deafness benefit
largely disappears ($-0.01$ to $+0.10$) while wheelchair use and plain language
keep theirs. Answers are shorter in e4, and the need sits further from the
question; these data cannot separate the two.

\section{Discussion}
\label{sec:discussion}

Models serve stated needs. That is the good news, and it is measured here against
the right control. The less good news is in the form of the service. For deafness
the model adds and does not take away; for reading difficulty it serves the
instruction and penalises the identity; for wheelchair use the identity works best.
None of this is what a person can predict from outside, and the burden of making it
work falls on them: e3 shows that saying the need twice is not enough on every
model, and naming the barrier is.

For auditors the methodological point carries furthest. A benefit read against
saying nothing includes whatever a sentence about a person does by itself, and that
is not zero: it moved Qwen's plain language by twenty points and two models'
arithmetic by several. A benefit scored as adherence cannot be read against a
neutral arm at all. Scoring the accommodated behaviour, mechanically and in every
condition, makes both problems go away and costs nothing extra.

The instrument is also a ready column for a benchmark of how well systems serve
people who are not the average user, which this research programme intends to
build.

\section{Limitations}
\label{sec:limits}

Three open 7--8B models at 4-bit on one wrapper; frontier models format
differently, and the screen-reader null would likely not hold for them. The checks
are lexical and conservative: the text-route lexicon misses ``send them an email'',
which undercounts additions in every cell. Pool questions were written to break
their need, so base rates are not estimates of general use. Two errors of ours
were found and corrected during the work, the plain-language splitter and the
deafness item set in e1, and one of interpretation, the wheelchair barrier; each is
reported where it occurs. The legal framing we first intended was dropped when the
statute did not support it. We did not test turn-by-turn disclosure, the ``declines
to say'' level, or how humans judge the answers.

\section{Reproduction}
\label{sec:repro}

Every script, stimulus, preregistration, raw answer and analysis output is in the
repository. Each experiment's analysis runs from its CSVs alone, without a GPU,
and each has an offline validator that injects known effects into a stubbed model
and checks that the analysis recovers them (46 checks for e0, 26 for e1, 15 for e2,
28 for e3 and e4 together, 36 for the instrument).

\begin{thebibliography}{99}

\bibitem{p2} M. Podgortsev. \emph{Single Human-AI Fit: One Attribute, Seven
Channels, Two Directions.} 2026. Concept DOI 10.5281/zenodo.22364497.

\bibitem{p3} M. Podgortsev. \emph{Stacked Human-AI Fit.} 2026.

\bibitem{kamruzzaman2025} M. Kamruzzaman, G. L. Kim. \emph{The Impact of Disability
Disclosure on Fairness and Bias in LLM-Driven Candidate Selection.}
arXiv:2506.00256; FLAIRS 2025.

\bibitem{accesseval2025} S. Panda, A. Agarwal, H. L. Patel. \emph{AccessEval:
Benchmarking Disability Bias in Large Language Models.} arXiv:2509.22703, 2025.

\bibitem{johnson2026} J. Johnson, A. Lewis, J. Mankoff, O. Banner. \emph{``I Don't
Trust it, but I Use it'': Navigating Trust, Privacy, and Identity in Disabled
People's Use of Generative AI.} CHI 2026. doi:10.1145/3772318.3790652.

\bibitem{gupta2026nd} I. Gupta, P. Buryi. \emph{How Frontier LLMs Adapt to
Neurodivergence Context: A Measurement Framework for Surface vs. Structural Change
in System-Prompted Responses.} arXiv:2605.00113, 2026.

\bibitem{penguin2025} Y. Wu, E. Sun, K. Zhu, J. Lian, J. Hernandez-Orallo,
A. Caliskan, J. Wang. \emph{Personalized Safety in LLMs: A Benchmark and A
Planning-Based Agent Approach.} arXiv:2505.18882; NeurIPS 2025.

\bibitem{curate} L. Alberts, A. Lupu, B. Ellis, J. Foerster. \emph{CURATe: Benchmarking
Personalised Alignment of Conversational AI Assistants.} arXiv:2410.21159.

\bibitem{prefeval2025} S. Zhao, M. Hong, Y. Liu, D. Hazarika, K. Lin. \emph{Do LLMs
Recognize Your Preferences? Evaluating Personalized Preference Following in LLMs.}
arXiv:2502.09597; ICLR 2025.

\bibitem{perg2025} C. Okite, N. Deng, K. Bodipati, H. Hou, J. Chai, R. Mihalcea.
\emph{Benchmarking and Improving LLM Robustness for Personalized Generation.}
arXiv:2509.19358, 2025.

\bibitem{pfqabench2026} Z. Sun, Y. Zhan, C. Shen, W. Yu, X. Zhang, M. He, J. Xu.
\emph{When Personalization Misleads: Understanding and Mitigating Hallucinations in
Personalized LLMs.} arXiv:2601.11000, 2026.

\bibitem{akpinar2025} N.-J. Akpinar, C.-J. Lee, V. Murdock, P. Perona. \emph{Who's
Asking? Evaluating LLM Robustness to Inquiry Personas in Factual Question
Answering.} arXiv:2510.12925, 2025.

\bibitem{yan2026} Y. Yan, S. Wu, T. Le. \emph{Locating and Controlling Implicit
Personalization in Large Language Models.} arXiv:2608.11735, 2026.

\bibitem{taheri2026} Taheri, Tazike, Carrington, Bigham. \emph{When Disability
Disclosure Travels: Memory, Privacy, and Contextual Integrity in Conversational AI.}
arXiv:2609.22720, 2026.

\bibitem{lachenmaier2026} Lachenmaier, Bultmann, Zarriess. \emph{Talking to a
Know-It-All GPT or a Second-Guesser Claude? How Repair reveals unreliable Multi-Turn
Behavior in LLMs.} arXiv:2604.19245, 2026.

\bibitem{laban2025} P. Laban, H. Hayashi, Y. Zhou, J. Neville. \emph{LLMs Get Lost
In Multi-Turn Conversation.} arXiv:2505.06120; ICLR 2026.

\bibitem{zhang2026} L. Zhang, J. Liu, T. August. \emph{Re-Centering Humans in LLM
Personalization.} arXiv:2606.06614, 2026.

\bibitem{crpd} United Nations. \emph{Convention on the Rights of Persons with
Disabilities}, Article 2, Definitions. 2006.

\bibitem{eqa} \emph{Equality Act 2010} (UK), section 20, Duty to make adjustments.

\end{thebibliography}
\end{document}
